% Generated draft: requires mathematical and source review.
\documentclass[11pt]{article}
\usepackage[margin=25mm]{geometry}
\usepackage{fontspec}
\setmainfont{Latin Modern Roman}
\usepackage{amsmath,amssymb,mathtools,unicode-math}
\setmathfont{Latin Modern Math}
\usepackage{enumitem,needspace,xurl,hyperref}
\setlength{\parindent}{0pt}
\setlength{\parskip}{5pt}
\setlength{\emergencystretch}{3em}
\newcommand{\N}{\mathbb{N}}
\newcommand{\Z}{\mathbb{Z}}
\newcommand{\Q}{\mathbb{Q}}
\newcommand{\R}{\mathbb{R}}
\newcommand{\C}{\mathbb{C}}
\newcommand{\F}{\mathbb{F}}
\newcommand{\A}{\mathbb{A}}
\newcommand{\PP}{\mathbb{P}}
\newcommand{\E}{\mathbb{E}}
\newcommand{\eps}{\varepsilon}
\newcommand{\dd}{\,\mathrm{d}}
\begin{document}
\section*{UnsolvedMath ID: 30005374}
\textbf{Source problem code:} catalogue entry
\subsection{Definitions and mathematical statement}
The supplied record does not determine an identifiable conjecture. Its opening word ``Conjecture'' is followed by explanatory prose, a theorem attributed to Bryan, and an incomplete alternate proof of an unstated Theorem 1. The following is a conditional reconstruction of the contextual genus statement, not a recovered open problem.

Let \(M=\mathbb{CP}^{2}\#\mathbb{CP}^{2}\) have its standard positive definite orientation, and let \(e_1,e_2\in H_2(M;\Z)\) be the standard generators, with \(e_i\cdot e_j=\delta_{ij}\). Thus the class denoted by the corrupted source text ``p2, 6q'' is interpreted as \(2e_1+6e_2\). For a smooth, closed, connected, oriented surface \(\Sigma\) smoothly embedded in \(M\), write \([\Sigma]\) for its integral homology class and \(g(\Sigma)\) for its genus. Under this conventional interpretation, the contextual assertion labeled ``Theorem 5 (Bryan [3])'' is
\[
\forall\,\Sigma\subset M,\qquad [\Sigma]=2e_1+6e_2\ \Longrightarrow\ g(\Sigma)\geq 10.
\]
The source says only ``a smooth surface representing'' the class; closedness, connectedness, orientability, embeddedness, and the choice of homology basis are not explicitly specified in the extracted theorem sentence. They are stated here as interpretive hypotheses requiring review, rather than silently supplied as verified source wording.

The surrounding passage attributes an alternate argument to minimum-genus bounds from gauge theory, in place of Floer-theoretic techniques attributed to Dai--Kang--Mallick--Park--Stoffregen. It describes \(2e_1+6e_2\) as twice a characteristic class. Here a class \(c\in H_2(M;\Z)\) is characteristic when \(c\cdot x\equiv x\cdot x\pmod 2\) for every \(x\in H_2(M;\Z)\). The passage asserts that the double cover \(X\) branched along \(\Sigma\) is smooth and spin and records
\[
 b_2(X)=4+2g(\Sigma),\qquad \sigma(X)=-16,
\]
where \(b_2\) is the second Betti number and \(\sigma\) is the signature of the intersection form. It attributes the signature computation to the \(G\)-signature theorem [2], and states that Furuta's ``\(10/8+2\)'' theorem [7] gives the intermediate bound \(g(\Sigma)\geq9\), improved to \(g(\Sigma)\geq10\) by Bryan's analysis of the covering involution. These are supplied contextual claims, not an independently checked proof; the full hypotheses of the cited results and of the branched-cover calculations are unavailable here.

The incomplete application describes a properly embedded annulus in \(M\) with the interiors of two disjoint \(4\)-balls removed, viewed as a cobordism from \(S^3\) to itself obtained from \(S^3\times I\) by two \(+1\)-framed \(2\)-handle attachments along an unlink. One annular boundary component is identified as the \((2,1)\)-cable of the figure-eight knot. Each attaching curve is said to link the annulus geometrically six times; with the indicated orientations, their algebraic linking numbers are six and two. The referenced figure is absent, cable conventions are unspecified, and the passage ends before describing the effect of the attachments. Consequently neither the full annulus construction nor the statement of Theorem 1 can be recovered.

The registry supplies the status label ``open'', but also marks the statement ``unrecoverable'' and supplies a verification note identifying this extraction as prose surrounding a proved argument. No precise open conjecture is asserted by this draft.
\subsection{Short English statement}
For permanent ID 30005374, the intended open problem cannot be recovered. The identifiable contextual theorem, attributed to Bryan, says that a smooth surface representing \((2,6)\) in \(\mathbb{CP}^{2}\#\mathbb{CP}^{2}\) has genus at least \(10\), subject to the explicitly stated interpretive hypotheses. The neighboring cable argument is incomplete. The registry's ``open'' label is supplied curation, not a verified mathematical conclusion.
\subsection{Sources}
\begin{itemize}
\item Local UnsolvedMath registry, supplied complete record: permanent ID 30005374; exact source problem code \texttt{catalogue entry}; title ``Background on Cable Sliceness and Genus Bounds''; Oberwolfach Reports Open Problems collection; proposed year 2023. No proposer is identified.
\item Supplied report citation: \emph{Low-dimensional Topology} (2023), Oberwolfach Reports, pp.\ 151--213, DOI 10.4171/owr/2023/3. Source URL and supplied literature-source link labeled ``original Oberwolfach report'': \url{https://doi.org/10.4171/owr/2023/3}. This authoritative report citation is supplied, but has not been independently consulted for this draft.
\item The extract attributes Theorem 5 and the improved genus bound to Bryan [3], the ``\(10/8+2\)'' theorem to Furuta [7], the signature calculation to the \(G\)-signature theorem [2], and the alternative Floer-theoretic techniques to Dai--Kang--Mallick--Park--Stoffregen. Complete bibliographic entries for [2], [3], and [7], and the identity of the speaker giving the alternate argument, are not supplied. The passage says that Furuta's preprint motivated Bryan's work.
\item Literature assessment checked at \texttt{2026-08-22}: supplied dated curation, not independently verified. It labels the record ``open'' because the extracted conjecture is not identifiable and describes the surrounding material as a proof of a neighboring cable non-sliceness theorem. The background explicitly requests independent verification of time-sensitive status and cited claims.
\item The supplied statement-verification note says that the official report confirms this extraction is explanatory prose surrounding a proved argument rather than a separately stated open problem. That verification is a supplied registry assertion; this draft does not claim to have visited the report. The record retains its original passage for provenance and supplies no clean statement.
\end{itemize}
\end{document}
